\label{chap:83-11}

\section{Identidad exponencial del TRIT y coordinación de los regímenes elíptico, parabólico e hiperbólico}

\subsection{Construcción del operador elíptico mediante la matriz de Paley y la relación \texorpdfstring{\(J_{+1}^{2}=-I\)}{J² = -I}}

La raíz de \(-1\) no se introduce como símbolo formal sin antecedente. Las
seis unidades módulo nueve se coordinatizan sobre
\(\mathbb P^1(\mathbb F_5)\), y el carácter cuadrático de
\(\mathbb F_5\) determina la matriz de conferencia de Paley. Su reducción
ternaria es
\[
 A_W=
 \begin{pmatrix}
 0&1&1&1&1&1\\
 1&0&1&2&2&1\\
 1&1&0&1&2&2\\
 1&2&1&0&1&2\\
 1&2&2&1&0&1\\
 1&1&2&2&1&0
 \end{pmatrix}\in M_6(\mathbb F_3).
\]
La identidad de conferencia reduce a
\begin{equation}\label{p12-83-c11-s1:eq:constantes-raiz-interna-menos-uno}
 \boxed{A_W^2=2I_6=-I_6.}
\end{equation}
Por tanto \(A_W\) es un operador de cuarto de giro en el espacio ternario.
Después de extender escala a \(\mathbb F_9\), se descompone en los espacios
propios conjugados de las raíces de \(X^2+1\). La fuente integral común
\[
 \mathcal Q_{\mathbb Z}=\mathbb Z[u]/(u^2+1)
\]
posee una realización finita \(u\mapsto A_W\) y una realización real
\(u\mapsto J_{\rm e}\). La fuente coordina la relación cuadrática; no
identifica cuerpos de características distintas.

\subsection{Familia unificada \texorpdfstring{\(\exp(tJ_\tau)=C_\tau(t)I+S_\tau(t)J_\tau\)}{exp(t J tau) = C tau(t) I + S tau(t) J tau}}

Los tres regímenes de TRIT se realizan mediante operadores
\[
 J_{+1}^2=-I,\qquad J_0^2=0,\qquad J_{-1}^2=I.
\]
Separando las potencias pares e impares de la exponencial se obtiene una
sola familia:
\begin{equation}\label{p12-83-c11-s1:eq:constantes-euler-extendido}
 \boxed{
 \exp(tJ_\tau)=C_\tau(t)I+S_\tau(t)J_\tau,
 \qquad \tau\in\{+1,0,-1\},}
\end{equation}
donde
\[
 (C_{+1},S_{+1})=(\cos,\sin),\qquad
 (C_0,S_0)=(1,t),\qquad
 (C_{-1},S_{-1})=(\cosh,\sinh).
\]
\subsection{Especializaciones elíptica, parabólica e hiperbólica y compatibilidad con el lector temporal}

Sus tres especializaciones rectoras son
\begin{equation}\label{p12-83-c11-s1:eq:constantes-euler-tres-ramas}
 \exp(\pi J_{+1})+I=0,\qquad
 \exp(J_0)=I+J_0,\qquad
 \exp\!\bigl(2(\log\varphi)J_{-1}\bigr)=F_{\rm av}^{\,2}.
\end{equation}
La rama elíptica contiene el cierre circular complejo; la parabólica, el
transporte de umbral; la hiperbólica, la autoescala áurea. La ecuación de
Euler aparece así como una especialización de una familia que coordina giro,
frontera y expansión hiperbólica.


\section{Fibración de Hopf, círculos de Villarceau y holonomía de Möbius como realización topológica del régimen elíptico}

\subsection{Estructura compleja y fibración de Hopf}

Sea \(S^3\subset\mathbb C^2\) y
\[
h(z_0,z_1)=
\bigl(2\Re(z_0\overline z_1),
2\Im(z_0\overline z_1),|z_0|^2-|z_1|^2\bigr)
\in S^2.
\]
En el toro de Hopf
\[
T_\eta=\{(z_0,z_1):|z_0|=\cos\eta,
|z_1|=\sin\eta\},
\]
las curvas
\[
\gamma_+(t)=(\cos\eta\,e^{it},\sin\eta\,e^{it}),
\]
\[
\gamma_-(t)=(\cos\eta\,e^{it},\sin\eta\,e^{-it})
\]
poseen comportamientos distintos: \(\gamma_+\) es una fibra de Hopf y
\(\gamma_-\) se proyecta con grado dos. Bajo proyección estereográfica, las
clases homológicas \((1,1)\) y \((1,-1)\) producen las dos familias diagonales
de Villarceau.

La conjugación compleja total
\[
K(z_0,z_1)=(\overline z_0,\overline z_1)
\]
preserva cada familia no orientada y revierte su orientación. No intercambia
las dos familias. La conjugación parcial puede intercambiarlas, pero no es una
aplicación compleja lineal ni el antiunitario estándar global. Esta corrección
es esencial para no fabricar una simetría partícula--antipartícula a partir del
dibujo toroidal.

\subsection{Doble recubrimiento y retorno spinorial}

La aplicación \(SU(2)\to SO(3)\) es un doble recubrimiento. Una vuelta de
\(2\pi\) en \(SO(3)\) puede levantar a \(-I\) en \(SU(2)\); la vuelta de
\(4\pi\) recupera \(I\). La estructura de signo no convierte un círculo en
banda de Möbius. Para obtener una banda se necesita un fibrado real de línea
cuyo carácter de holonomía sea \(-1\).

Sea \(L_\chi\to S^1\) el fibrado asociado a
\[
\chi:\pi_1(S^1)\longrightarrow\{\pm1\},
\qquad \chi(1)=-1.
\]
Su espacio total es una banda de Möbius. La ruta \(C_M=-\operatorname{rev}\)
aporta el candidato de signo; el círculo Villarceau aporta el lazo; el
entrelazador físico debe probar que el primero es la holonomía del segundo.

\subsection{Correspondencia entre el espectro \texorpdfstring{\(\{7,3\}\)}{{7,3}}, la razón \texorpdfstring{\(5{:}2\)}{5:2} y el retorno spinorial}

Con \(B_c=7I+2S\), defínase el operador orientado de memoria
\[
B_{\rm mem}^{\varepsilon}=5I+2\varepsilon S,
\qquad\varepsilon=\pm1.
\]
Para \(\varepsilon=+1\),
\[
\Spec(B_{\rm mem}^{+})=\{7,3\},
\qquad\det B_{\rm mem}^{+}=21.
\]
Interpretar los valores \((7,3)\) como bordes ecuatoriales exterior e interior
en una carta positiva común produce
\[
R+r=7\ell,\qquad R-r=3\ell,
\]
y por tanto
\[
R=5\ell,\qquad r=2\ell,\qquad
\sin\beta=\frac rR=\frac25.
\]
El ángulo Villarceau es
\[
\beta=\arcsin\frac25
=23.5781784782\ldots^\circ.
\]

El teorema algebraico \(\{7,3\}\leftrightarrow(5,2)\) es exacto. La premisa
geométrica que interpreta el espectro como los dos bordes de un toro es una
interfaz declarada. No debe ocultarse ni emplearse el ángulo como blanco.

La realización del bloque completo \(B_c\) conserva otra lectura:
\(R:r=7:2\) y \(\beta=\arcsin(2/7)\). No son valores rivales; pertenecen a
dos funtores: bloque total y modo de memoria. La prioridad física del segundo
debe decidirse por el entrelazador con la dinámica, no por cercanía decimal con
otra constante.

\subsection{Representación del par angular en el régimen elíptico y cierre de la unidad de fase}

El ángulo nativo de HMT es una fase \(u\in\mathbb R/\mathbb Z\). Los lectores
publican
\[
\theta_{\rm rad}=2\pi u,
\qquad
\theta_{\rm deg}=360u.
\]
Radianes y grados son cartas distintas del mismo elemento circular. El radián
no es más «ontológico» por ser adimensional; su naturalidad procede de la
longitud de arco. Los grados proporcionan otra coordenada absoluta de la
vuelta. Cualquier relación con \(\alpha^{-1}\), \(\varphi\) o CKM debe
formularse primero como cociente o holonomía invariante y sólo después
evaluarse en una carta angular.

La identidad exacta ya disponible
\[
\tanh^2(2\log\varphi)=\frac59
\]
conecta el canal Perron--Fibonacci con el cociente espectral de \(B_c\). Su
complemento es \(4/9\), el salto normalizado. Esta relación es estructural.
No implica que todo ángulo construido con \(\varphi\) sea un ángulo físico.

\subsection{Brouwer, Möbius y grado}

El teorema del punto fijo de Brouwer, las transformaciones fraccionarias de
Möbius y la banda de Möbius son tres objetos distintos. Se relacionan sólo
después de fijar dominio, acción y fibrado. Las órbitas elípticas de adelanto y
retardo pueden modelarse como dos secciones de un fibrado con holonomía de
signo; la transformación de Möbius describe la acción conformal en la carta;
el teorema de Brouwer garantiza un punto fijo para mapas continuos del disco.
Ninguna de estas frases sustituye el diagrama de realización.

\begin{fronteraBox}[title={Resultado y frontera}]
Se demuestran el grado doble de la curva Hopf, la orientación de la
conjugación total y la construcción algebraica \(7,3\mapsto5,2\). La lectura
toroidal (5{:}2), la holonomía Möbius y la interpretación
partícula--antipartícula forman una cadena de interfaces que debe compartir un
único entrelazador. El volumen no fuerza ese entrelazador mediante una
coincidencia angular.
\end{fronteraBox}
